% Complete finite proofs for the first structural causal model component.
\section{Finite structural causal models: proofs}
\label{sec:annex-p16-finite-causal-models}

Use the graph, probability and equation-replacement definitions in
\cref{sec:p16-finite-causal-models}. The source supplies the structural
interpretation, intervention operation, identification criterion and
factorization statements
\cite[secs.~3.1--3.2.3, pp.~103--112]{Pearl2009CausalOverview}.
The finite constructions and proofs here include explicit treatment of
zero probabilities and dependent disturbances.

\begin{annexitem}{annex:p16.causal-topological-order}{thm:p16-causal-topological-order}{Ordering an acyclic finite graph}
\begin{proof}
First consider a nonempty finite acyclic directed graph. Suppose that every
vertex had an incoming edge. Choose a vertex \(j_0\), then choose an
incoming predecessor \(j_1\to j_0\), then a predecessor
\(j_2\to j_1\), and continue until one more vertex has been selected
than the graph contains. Some vertex repeats, say \(j_a=j_b\)
with \(a<b\). The edges
\(j_b\to j_{b-1}\to\cdots\to j_a\) form a directed cycle,
contrary to acyclicity. Therefore some vertex has no incoming edge.

Remove such a vertex and all its incident edges. The remaining graph is
acyclic, because any cycle in it would be a cycle in the original graph.
Whenever the remaining graph is nonempty, the preceding argument supplies
another vertex to remove. Exactly one distinct vertex is removed at each
step, so the procedure stops after \(n\) steps. Let
\(i_1,\ldots,i_n\) be the removal order. If \(j\to i_r\), then
\(j\) cannot still be in the graph when \(i_r\) is removed, because
it would supply an incoming edge. Thus \(j\) occurs earlier in the
order. This proves the topological property and termination.
\end{proof}
\end{annexitem}

\begin{annexitem}{annex:p16.causal-solution}{thm:p16-causal-solution}{Unique solutions and probability laws}
\begin{proof}
Fix \(\omega\in\Omega\) and a topological order
\(i_1,\ldots,i_n\). At the first vertex the parent tuple is empty.
The total function \(f_{i_1}\), applied to that tuple and
\(U_{i_1}(\omega)\), gives a member of \(S_{i_1}\). At step
\(r\), every parent of \(i_r\) has already been assigned a value,
so the total function \(f_{i_r}\) gives its value. Subsequent steps
never change earlier coordinates. Consequently every structural equation
is satisfied by the completed tuple.

If \(v,w\) are two solutions at \(\omega\), their first
coordinates in the topological order agree, since each applies the same
function to the same empty parent tuple and the same disturbance. If their
earlier coordinates agree, the parents of \(i_r\) agree and the
equation for \(i_r\) forces its values to agree. Induction gives
\(v=w\). Any other topological evaluation is also a solution, hence
must give this same tuple. Since this holds separately at every outcome,
it defines functions \(V_i:\Omega\to S_i\).

For each \(v\in S_N\), the set \(\{\omega:V(\omega)=v\}\)
is an event in the finite probability space. These sets are pairwise
disjoint and their union is \(\Omega\); empty fibers contribute zero.
Thus each stated mass is nonnegative and
\[
 \sum_{v\in S_N}p_M(v)
 =\sum_{v\in S_N}\sum_{\omega:V(\omega)=v}q(\omega)
 =\sum_{\omega\in\Omega}q(\omega)=1.
\]
For a fixed \(J\subseteq N\) and \(s\in S_J\), partition the
event \(\{V_J=s\}\) by its full assignments \(v\) with
\(v_J=s\). Finite additivity gives the marginal formula. This
argument also covers \(J=\varnothing\): there is just one empty
tuple, of probability one. No division by a probability is involved.
\end{proof}
\end{annexitem}

\begin{annexitem}{annex:p16.causal-intervention-laws}{thm:p16-causal-intervention-laws}{Equation replacement, consistency and unaffected vertices}
\begin{proof}
Deleting incoming edges cannot introduce a directed cycle, so the
intervention graph is again a DAG. The replacement functions are total
constant functions into their original value sets. All other functions
and the finite probability space remain valid model data. The preceding
solution theorem therefore applies to \(M_{I=x}\), proving its
unique solution and the mass formula. The equations at vertices in
\(I\) force its solution to have coordinates \(x\), so assignments
with different coordinates have empty fibers and probability zero.

For the first pointwise assertion, suppose \(V_I(\omega)=x\).
The original tuple \(V(\omega)\) already satisfies every unchanged
equation, and by the supposition it also satisfies every replacement
equation. It is therefore a solution of the modified system at
\(\omega\). Uniqueness implies
\(V^{I=x}(\omega)=V(\omega)\).

For the second assertion, put
\(H=N\setminus(I\cup\operatorname{de}_G(I))\). Every parent of
a vertex in \(H\) is in \(H\): a parent in \(I\) gives a
length-one path from \(I\), and a parent in
\(\operatorname{de}_G(I)\) gives such a path after appending its
edge. Either possibility would contradict membership of the child in
\(H\). Read the original topological order restricted to \(H\).
At its first vertex the parent inputs, if any, cannot occur outside
\(H\) or later in the order. At each vertex in \(H\), its equation
and disturbance are unchanged, and its parent values agree by induction.
The value at that vertex therefore agrees in the two evaluations. This
proves the pointwise assertion throughout \(H\).

The empty intervention replaces no equation and deletes no edge. For
disjoint \(I,J\), each original equation belongs to exactly one of
three cases: it is replaced by the prescribed constant in \(I\),
replaced by the prescribed constant in \(J\), or left unchanged.
The order of these two replacements does not change any case, graph,
probability weight or disturbance function. The resulting models are the
same, so their solutions and distributions are the same as well.
\end{proof}
\end{annexitem}

\begin{annexitem}{annex:p16.causal-factorization}{thm:p16-causal-factorization}{Independent disturbances give the structural product}
\begin{proof}
We first record the finite consequence of the stated mutual independence.
For arbitrary subsets \(A_i\subseteq\mathcal U_i\), partition the
event \(\bigcap_i\{U_i\in A_i\}\) by its disturbance tuples.
Then finite distributivity gives
\[
\begin{aligned}
 \mathbb P\Bigl(\bigcap_{i\in N}\{U_i\in A_i\}\Bigr)
 &=\sum_{u_1\in A_1}\cdots\sum_{u_n\in A_n}
                         \prod_{i\in N}r_i(u_i)\\
 &=\prod_{i\in N}\sum_{u_i\in A_i}r_i(u_i).
\end{aligned}
\]
An empty set gives zero on both sides. Taking \(A_i=\mathcal U_i\)
for omitted indices proves the same identity for a subset of the
disturbances, because those factors sum to one.

Fix \(v\in S_N\), and define
\[
 A_i(v)=\{u\in\mathcal U_i:
               f_i(v_{\operatorname{pa}_G(i)},u)=v_i\}.
\]
For each \(\omega\), the assertion \(U_i(\omega)\in A_i(v)\)
for every \(i\) says exactly that the fixed tuple \(v\) satisfies
all structural equations at that outcome. By uniqueness this holds
exactly when \(V(\omega)=v\). Therefore
\[
 p_M(v)=\mathbb P\Bigl(\bigcap_i\{U_i\in A_i(v)\}\Bigr)
       =\prod_i\sum_{u\in A_i(v)}r_i(u)
       =\prod_i k_i(v_i\mid v_{\operatorname{pa}_G(i)}).
\]
This proves the product at every tuple, without conditioning.

For the conditional assertion, take a topological order in which vertex
\(i\) occurs in position \(r\), and let \(L\) be the set of
vertices in earlier positions. Recursive evaluation shows that the tuple
\(V_L\) is a function of \((U_j)_{j\in L}\): at each earlier
step its parents also lie earlier, so substituting their already
constructed functions introduces no disturbance outside \(L\).
In particular \(V_{\operatorname{pa}_G(i)}\) is a function of
this disturbance tuple, and \(i\notin L\).

For a tuple \(u_L\), summing the joint disturbance mass over all
indices outside \(L\cup\{i\}\) gives
\[
 \mathbb P(U_i=u,(U_j)_{j\in L}=u_L)
       =r_i(u)\prod_{j\in L}r_j((u_L)_j).
\]
When \(L\) is empty, its tuple has probability one and the same
identity holds. Sum this identity over the tuples \(u_L\) whose
evaluation gives a specified parent assignment \(a\). These tuple
events partition \(\{V_{\operatorname{pa}_G(i)}=a\}\). Hence
\[
 \mathbb P(U_i=u,V_{\operatorname{pa}_G(i)}=a)
       =r_i(u)\mathbb P(V_{\operatorname{pa}_G(i)}=a).
\]
On this parent event, the equation for \(i\) gives
\(V_i=f_i(a,U_i)\). Summing the last identity over the disturbance
values with output \(b\) yields
\[
 \mathbb P(V_i=b,V_{\operatorname{pa}_G(i)}=a)
       =k_i(b\mid a)\mathbb P(V_{\operatorname{pa}_G(i)}=a).
\]
Division by the assumed positive final factor proves the conditional
identity. If that factor is zero, the displayed joint identity still
holds, but it supplies no conditional ratio. The structural kernel
remains defined by its original finite sum.
\end{proof}
\end{annexitem}

\begin{annexitem}{annex:p16.causal-truncation}{thm:p16-causal-truncation}{Truncating the product after replacement}
\begin{proof}
The modified model has the same disturbance functions and their same
joint law, so mutual independence still holds. For \(i\notin I\)
its function and parent set are unchanged, so its structural kernel is
the original \(k_i\). For \(i\in I\) its parent set is empty
and its function has constant value \(x_i\). Its new kernel is
\[
 \sum_{u\in\mathcal U_i}\mathbf1_{\{b=x_i\}}r_i(u)
       =\mathbf1_{\{b=x_i\}},
\]
since \(\sum_u r_i(u)=1\). Apply
\cref{thm:p16-causal-factorization} to this acyclic modified model.
The product contains one such indicator for each intervened coordinate
and the original factors at all other coordinates. If any intervened
coordinate disagrees with \(x\), its indicator makes the product zero.
If all agree, the indicators are one and the remaining product is exactly
the claimed expression, with those parent values fixed to \(x\).

The marginal formula is the grouping-of-fibers result already proved for
finite solution laws. When \(I=N\), every coordinate is fixed; the
remaining product is the empty product one, and the law has mass one
exactly at the single tuple \(x\). The proof uses the structural
kernels and never assigns a conditional ratio to a null event.
\end{proof}
\end{annexitem}

\begin{annexitem}{annex:p16.causal-positive-identification}{thm:p16-causal-positive-identification}{A positive observational law determines the kernels}
\begin{proof}
Let \(a\in S_{\operatorname{pa}_G(i)}\). All value sets are
nonempty and finite, so choose arbitrary values for the coordinates
outside the parent set to extend \(a\) to some \(w\in S_N\).
Full positive support gives \(p_M(w)>0\). Thus
\[
 \mathbb P(V_{\operatorname{pa}_G(i)}=a)
   =\sum_{z:z_{\operatorname{pa}_G(i)}=a}p_M(z)
   \geq p_M(w)>0.
\]
When the parent set is empty, the same probability is one. By
\cref{thm:p16-causal-factorization}, every kernel value is its
observational conditional ratio. Express its joint-event numerator and
parent-event denominator by their finite marginal sums. Substitution
into \cref{thm:p16-causal-truncation} gives the displayed product,
including the empty product when all coordinates are intervened upon.

Suppose now that two models in the specified class have equal full
observational laws. They have the same graph and value sets, so the
numerators and denominators in these ratios use the same index sets and
equal masses. They consequently have equal kernels and equal truncated
products at every assignment. Outside the intervention assignment both
laws are zero. This is equality of their full intervention laws and proves
identification by \cref{def:p16-causal-identification}.
Summing equal laws gives equal marginals; multiplying their masses by a
fixed finite response function and summing gives equal response means.
No inference about estimation from a finite sample enters this argument.
\end{proof}
\end{annexitem}

\begin{annexitem}{annex:p16.causal-confounded-example}{thm:p16-causal-confounded-example}{Computing the common-cause model}
\begin{proof}
The proposed outcome weights sum to
\[
 \sum_{z=0}^1\sum_{a=0}^1\sum_{b=0}^1
          \tfrac12\alpha(a)\alpha(b)
 =\Bigl(\sum_{z=0}^1\tfrac12\Bigr)
    \Bigl(\sum_{a=0}^1\alpha(a)\Bigr)
    \Bigl(\sum_{b=0}^1\alpha(b)\Bigr)=1.
\]
Their factorization gives the stated independent coordinate disturbances.
The graph's order \(Z,X,Y\) evaluates to
\(V(z,a,b)=(z,z\oplus a,z\oplus b)\).
For given \(z,x,y\), the unique disturbance coordinates giving
that assignment are \(a=x\oplus z\), \(b=y\oplus z\).
Consequently
\[
 p_{M_A}(z,x,y)=\tfrac12\alpha(x\oplus z)\alpha(y\oplus z).
\]
For \(x=y=0\), summing over \(z\) gives
\[
 \tfrac12\bigl((\tfrac34)^2+(\tfrac14)^2\bigr)=\tfrac5{16}.
\]
For \(x=y=1\) the two terms are reversed, giving the same value.
For either unequal pair, the sum is
\[
 \tfrac12\bigl((\tfrac34)(\tfrac14)
                    +(\tfrac14)(\tfrac34)\bigr)=\tfrac3{16}.
\]
This proves every table entry. Both rows and both columns sum to
\(8/16=1/2\). Dividing each \(Y=1\) entry by its positive row
sum gives \(5/8\) for \(X=1\) and \(3/8\) for \(X=0\).

Under \(\operatorname{do}(X=x)\), the equations instead evaluate
to \((Z,X,Y)=(z,x,z\oplus b)\). The outcome weights remain
\(\tfrac12\alpha(a)\alpha(b)\). The \(Y=1\) event consists
of \(z=0,b=1\) or \(z=1,b=0\), with arbitrary \(a\).
Its mass is
\[
 \tfrac12\alpha(1)\sum_a\alpha(a)
 +\tfrac12\alpha(0)\sum_a\alpha(a)
 =\tfrac18+\tfrac38=\tfrac12.
\]
This expression does not depend on \(x\). For a bit-valued response,
its expectation equals its probability of being one, by summing the
zero and one values. The observational difference is therefore
\(5/8-3/8=1/4\), and the intervention difference is
\(1/2-1/2=0\).

For the adjustment identity, sum the three-variable law over \(y\):
\[
 \mathbb P(X=x,Z=z)=\tfrac12\alpha(x\oplus z)>0.
\]
The numerator for \(Y=1\) is this same positive factor multiplied
by \(\alpha(1\oplus z)\). Dividing yields
\[
 \mathbb P(Y=1\mid X=x,Z=z)=\alpha(1\oplus z).
\]
It is \(1/4\) at \(z=0\) and \(3/4\) at \(z=1\).
Weighting by \(\mathbb P(Z=z)=1/2\) gives
\(1/8+3/8=1/2\), the intervention probability already calculated.
Every division here has an explicit positive denominator. This is a
calculation in the displayed model, with no appeal to a general graphical
adjustment theorem.
\end{proof}
\end{annexitem}

\begin{annexitem}{annex:p16.causal-nonidentification}{thm:p16-causal-nonidentification}{Two pairs of observationally equal models}
\begin{proof}
For the first assertion, define \(M_A'\) with endogenous variables
only \(X,Y\), each binary, and allowed graph \(X\to Y\).
Use the same three-bit outcome space and weights as \(M_A\).
Give its two disturbances value sets \(\{0,1\}^2\), with
\[
 \widetilde U_X(z,a,b)=(z,a),\qquad
 \widetilde U_Y(z,a,b)=(z,b).
\]
Define the structural functions by
\[
 f_X((),(z,a))=z\oplus a,\qquad
 f_Y(x,(z,b))=z\oplus b.
\]
Here \(()\) denotes the empty parent tuple. The second function is
total on its parent input \(x\) and disturbance, and ignores \(x\).
This is permitted by the model definition. Its graph is acyclic; its
disturbances need not be independent in the stated class. Evaluation
gives exactly the \(X,Y\) projection of \(M_A\), both before and
after replacing the equation for \(X\). Thus its observational law
is the displayed table and its target is \(1/2\).

For comparison, define \(M_B\) on outcome space \(\{0,1\}^2\)
with coordinates \((c,d)\) and weights
\[
 q_B(c,d)=\tfrac12\beta(d),\qquad
 \beta(0)=\tfrac58,\quad\beta(1)=\tfrac38.
\]
These weights are nonnegative and sum to
\((1/2+1/2)(5/8+3/8)=1\). Give it coordinate disturbances
\(U_X=c\), \(U_Y=d\), the same binary endogenous value sets,
the same graph \(X\to Y\), and equations
\[
 X=U_X,\qquad Y=X\oplus U_Y.
\]
The disturbances are independent; hence this model also belongs to the
larger class allowing dependence. Its solution is \((c,c\oplus d)\).
For a fixed observed pair \((x,y)\), the sole compatible outcome
is \((c,d)=(x,x\oplus y)\). Therefore
\[
 p_{M_B}(x,y)=\tfrac12\beta(x\oplus y)
  =\begin{cases}5/16,&x=y,\\3/16,&x\ne y.\end{cases}
\]
This agrees at all four pairs with \(p_{M_A'}\), and all four
masses are positive. Under \(\operatorname{do}(X=1)\), however,
\(Y=1\oplus d\), which is one exactly when \(d=0\). Its
probability is \(5/8\), unequal to \(1/2\). The two models
have the same specified graph, observed variables and observed law, but
different targets. They are the required nonidentification witness.

For the second assertion, use the same graph and binary endogenous value
sets, but let the outcome space be a singleton of mass one. Both
disturbance value sets are singletons and both disturbances are constant.
They are mutually independent: the sole joint tuple and each sole
marginal value have probability one. In \(M_C\) put \(X=0\)
and \(Y=0\), with the function for \(Y\) ignoring its allowed
parent. In \(M_D\) put \(X=0\) and \(Y=X\).
All functions are total on their full finite parent domains.
Both models evaluate observationally to \((0,0)\), so their laws
are equal point masses. Under \(\operatorname{do}(X=1)\), the
first gives \(Y=0\) and the second gives \(Y=1\). Their targets
are zero and one. The observational conditional on \(X=1\) is
undefined in both models, and the observation law supplies no value for
the response kernel at that parent input. This proves the claimed
nonidentification without the positive-support restriction.
\end{proof}
\end{annexitem}
